\chapter{Bone Marrow Tests}

\section*{What Is This Test?}
Bone marrow tests examine marrow cells and tissue to evaluate blood-cell production, leukemia, lymphoma, myeloma, unexplained cytopenias, marrow failure, and some infections or storage disorders. The procedure commonly includes aspiration and a core biopsy.

\section*{Alternative Names}
\begin{itemize}
  \item Bone Marrow Aspiration
  \item Bone Marrow Biopsy
  \item Marrow Examination
  \item Trephine Biopsy
\end{itemize}

\section*{Principle}
An aspirate provides individual cells for morphology, flow cytometry, cytogenetics, and molecular testing. A core biopsy preserves tissue architecture and estimates cellularity, fibrosis, infiltration, and focal disease.

\section*{Why This Test Is Ordered}
It is used for persistent unexplained anemia, leukopenia, leukocytosis, thrombocytopenia, pancytopenia, abnormal blood smear, suspected hematologic cancer, staging or monitoring selected cancers, and evaluation of treatment response.

\section*{Primary vs Secondary Test}
Bone marrow examination is a definitive or disease-classifying test in many hematologic evaluations. CBC, smear review, flow cytometry, imaging, and molecular tests help determine whether it is needed and how the specimen should be studied.

\section*{How This Test Is Performed}
After local anesthesia, a needle is usually placed into the back of the pelvic bone. Liquid marrow is aspirated and a small core is removed. The samples are sent for morphology and any ancillary studies requested by the clinical question.

\section*{What Preparation Is Needed}
Tell the team about bleeding disorders, low platelets, anticoagulants, allergies, pregnancy possibility, and prior marrow results. Sedation may be offered; arrange transportation if it is used.

\section*{Contraindications and Precautions}
Temporary pain, bruising, bleeding, and infection are possible. Medication changes must be individualized. A dry tap can occur with fibrosis or tightly packed marrow and may require reliance on the core biopsy or another sampling approach.

\section*{Understanding Results}
Results may describe cell lines, blasts, dysplasia, abnormal infiltrates, fibrosis, iron stores, chromosomes, or molecular changes. A diagnosis depends on morphology, immunophenotype, genetics, blood findings, symptoms, and imaging rather than one result alone.

\section*{Follow-up Tests}
Follow-up may include repeat CBC and smear, flow cytometry, karyotype or FISH, molecular studies, lymph-node or other tissue biopsy, imaging, cerebrospinal-fluid assessment, or response monitoring.

\section*{References}
\begin{enumerate}
  \item MedlinePlus. Bone Marrow Tests.
  \item College of American Pathologists. Bone marrow examination resources.
  \item World Health Organization. Classification of Haematolymphoid Tumours.
\end{enumerate}

\clearpage
\section*{Understanding Results and Follow-up}
The laboratory report should identify the specimen, method, units, reference interval or decision limit, and whether the result is positive, negative, abnormal, or inconclusive. Reference intervals vary with age, sex, pregnancy, specimen type, assay, and laboratory; a result outside a range is not automatically a diagnosis, and a result within a range does not always exclude disease. Discuss unexpected results with the ordering clinician, who can decide whether repeat or confirmatory testing, treatment, or referral is appropriate. Do not change medicines or delay urgent care based on an isolated result.


\section*{Regulatory Status and Authorization Date}
This chapter describes a test category, not a specific commercial product. FDA, CDSCO, EU IVDR, and MHRA approval or registration dates therefore cannot be assigned without the assay manufacturer, product name, instrument, and intended use. For laboratory-developed tests, an individual agency approval date may not exist. See the Preface for the interpretation of regulatory status.

\clearpage
